\chapter{完整回归结果}\label{app:regress}

\section{完整 RTL 回归}
2026 年 10 月 14 日完整回归共 \RegCount{} 项，全部通过（表~\ref{tab:regress}）。自检合计 \RegChecks{} 次，独立核对合计 \RegIndep{} 项；“—”表示未接入参考检查器。测试平台下方为非默认参数。

\begin{longtable}{@{}p{1.45cm}p{3.3cm}p{5.9cm}rr@{}}
\caption{完整 RTL 回归结果（均为通过）}\label{tab:regress}\\
\longtablehead{5}{编号 & 测试平台 & 测试内容 & 自检 & 独立核对}
BAS-01 & \code{tb_ram_sdp} & 核对不同存储器配置的写入、同步读取和读使能 & 20764 & — \\
MAC-01 & \code{tb_mac_tx} & 核对以太网组帧、帧间隔和上游缺字节处理 & 264649 & — \\
MAC-02 & \code{tb_mac_rx} & 核对正常帧、错误分类、去校验值和帧边界 & 105303 & — \\
MAC-03 & \code{tb_mac_loop} & 发送端接回接收端，核对不同帧长与双网内容 & 1385254 & — \\
TX-01 & \code{tb_tx_basic} & 核对不同消息长度的分页存储、协议头与双网发送 & 662055 & 7605 \\
TX-02 & \code{tb_tx_basic}\par \code{TX_PAGE_AW=5}\par \code{TX_OFF_AW=11} & 核对不同消息长度的分页存储、协议头与双网发送 & 661635 & 7605 \\
TX-03 & \code{tb_tx_bag} & 核对带宽间隔、计数回绕和发送间隔变化 & 14054170 & 14591 \\
TX-04 & \code{tb_tx_multi} & 核对多条虚拟链路的队列、轮询与网络选择 & 14438575 & 56211 \\
TX-05 & \code{tb_tx_drop} & 核对禁用、队满、超长、无空闲页及失败后的页回收 & 30267256 & 4642 \\
TX-06 & \code{tb_tx_net} & 用固定输入核对发送内容和事件，供网表对照 & 219313 & 229 \\
RX-01 & \code{tb_rx_basic} & 核对接收映射、不同消息长度、用户数据与端口 & 9960369 & 123855 \\
RX-02 & \code{tb_rx_fault} & 损坏校验值、长度、报头和配置，核对拒绝原因 & 576923 & 890 \\
RX-03 & \code{tb_rx_icrm} & 核对序号规则、回绕和两网到达时间差边界 & 4635288 & 1274 \\
RX-04 & \code{tb_rx_load} & 大量虚拟链路接收时，核对序号隔离和冗余去重 & 6406189 & 19180 \\
RX-05 & \code{tb_rx_ovf_ring}\par \code{RX_RING_AW=9} & 缩小字节缓存，核对溢出回退和后续接收恢复 & 422837 & 99 \\
RX-06 & \code{tb_rx_ovf_dq}\par \code{RX_DQ_AW=1} & 缩小交付队列，核对队满丢帧和释放后恢复 & 373852 & 57 \\
RX-07 & \code{tb_rx_net} & 用固定正常及异常输入核对接收输出，供网表对照 & 374111 & 110 \\
RX-08 & \code{tb_rm_phase} & 跨粗计时边界并交换处理顺序，核对去重和老化 & 353130 & — \\
SYS-01 & \code{tb_sys_basic} & 两个完整端系统互发消息，核对交付与事件 & 5155148 & 2896 \\
SYS-02 & \code{tb_sys_fault} & 在线路上丢帧、重放及损坏字节，核对检测和恢复 & 47117674 & 8726 \\
SYS-03 & \code{tb_sys_delay} & 改变两网延迟，核对交付、去重和端到端时延 & 15481734 & 3763 \\
SYS-04 & \code{tb_sys_load} & 多链路同时传输，核对交付、发送间隔和配置隔离 & 216810445 & 116677 \\
SYS-05 & \code{tb_sys_cfg} & 核对配置写握手、读延迟和计数采样时刻 & 21167623 & 714342 \\
SYS-06 & \code{tb_sys_reset} & 传输途中复位，核对旧残帧不混入新交付 & 10991314 & 3682 \\
SYS-07 & \code{tb_sys_ovf}\par \code{RX_RING_AW=9} & 完整系统接收缓存溢出时，核对丢弃、计数和恢复 & 9681100 & 1488 \\
SYS-08 & \code{tb_sys_net} & 通过顶层执行配置、正常及异常传输，供网表对照 & 2943996 & 1265 \\
SYS-09 & \code{tb_elinx} & 由界面入口运行双端双网互发，核对字节、端口和去重 & 9238 & — \\
SVC-01 & \code{tb_sys_snmp} & 核对查询、团体名和编码边界，并比较配置口读值 & 37611 & 926 \\
SVC-02 & \code{tb_sys_snmp_settrap} & 核对设置、非法写入拒绝、启动和认证失败通知 & 8563 & 188 \\
SVC-03 & \code{tb_judge_snmp} & 解码查询、设置和通知，核对字段与配置口读值 & 699 & — \\
SVC-04 & \code{tb_sys_615a} & 核对加载卸载、文件内容、校验错误和丢包重试 & 44098 & 752 \\
SVC-05 & \code{tb_sys_rtc} & 核对空载及有流量时的时钟同步、秒边界和交付 & 10675857 & 1811 \\
SVC-06 & \code{tb_rtc_recovery} & 丢弃同步回应并交换使能顺序，核对重发恢复 & 28 & — \\
SVC-07 & \code{tb_rtc_originate} & 错误或过期的请求时间不得校时，合法回应仍能同步 & 34 & — \\
PERF-01 & \code{tb_sys_perf} & 测量吞吐、发送与接收时延、端到端时延和抖动 & 756575324 & 1212916 \\
PERF-02 & \code{tb_gbe_burst}\par \code{COUNT=77} & 同时提交多链路消息，测量排队时间并逐字节查重 & 643550 & — \\
PERF-03 & \code{tb_gbe_burst}\par \code{COUNT=81} & 同时提交多链路消息，测量排队时间并逐字节查重 & 673498 & — \\
CAP-01 & \code{tb_judge_capacity} & 配置全部收发链路，逐字节核对双网只交付一次 & 49924 & — \\
\end{longtable}

\section{评委入口}
2026 年 10 月 14 日运行，7 个用例全部通过，Error 和 Warning 均为 0，总耗时 626.779~s。该运行与完整回归相互独立，检查数不相加。

\begin{longtable}{@{}p{3.3cm}p{6.0cm}r c r@{}}
\caption{评委入口七个用例结果}\\
\longtablehead{5}{用例 & 测试平台与内容 & 自检项数 & 结果 & 耗时（s）}
\code{services}（SVC-08） & \code{tb_net_services}：用户收发、SNMP、RTC、加载与卸载 & \SvcChecks & PASS & 2.411 \\
\code{capacity}（CAP-01） & \code{tb_judge_capacity}：每端发送 128 条、配置接收 4096 条 VL & \CapChecks & PASS & 2.462 \\
\code{snmp}（SVC-03） & \code{tb_judge_snmp}：查询、设置、通知及配置读回 & \SnmpChecks & PASS & 5.645 \\
\code{burst77}（PERF-02） & \code{tb_gbe_burst}：77 条 VL 同时提交 & 643550 & PASS & 2.454 \\
\code{burst81}（PERF-03） & \code{tb_gbe_burst}：81 条 VL 同时提交 & 673498 & PASS & 2.569 \\
\code{rtc}（SVC-05） & \code{tb_sys_rtc}：空载及带载同步误差 & 10675857 & PASS & 57.788 \\
\code{performance}（PERF-01） & \code{tb_sys_perf}：吞吐、时延与发送抖动 & 756575324 & PASS & 552.375 \\
\end{longtable}

\chapter{复现命令}\label{app:repro}

\section{环境与行为仿真入口}
命令在仓库根目录执行，PowerShell 命令行末的反引号为续行符。工程路径不得含中文或空格。环境：Windows PowerShell 5.1、eLinx 3.0.7、ModelSim SE-64 2020.4；Python 脚本要求 3.10 及以上版本（已验证 3.12.7）。工具路径按表~\ref{tab:envvars} 设置，详见 \code{README.md}“软件准备”。

\begin{table}
\centering
\caption{复现所需工具路径}\label{tab:envvars}
\begin{tabularx}{\textwidth}{@{}LL@{}}
\toprule
变量 & 取值 \\
\midrule
\code{AFDX_PYTHON} & 完整的 \code{python.exe} 路径；不设时使用 PATH 中的 Python \\
\code{MODELSIM_HOME} & 含 \code{vsim.exe}、\code{vlog.exe}、\code{vlib.exe} 的目录；不设时从 PATH 查找 \\
\code{ELINX_HOME} & 含 \code{bin/shell/bin} 的安装根目录；不设时从 PATH 查找综合程序 \\
\bottomrule
\end{tabularx}
\end{table}

Python 依赖的安装命令如下；界面行为仿真和纯 ModelSim 入口不需要 Python。
\begin{small}
\begin{verbatim}
python -m pip install psutil pyasn1==0.4.8
\end{verbatim}
\end{small}

\textbf{界面行为仿真：}打开 \code{elinx/AFDX_EndSystem.epr}，在 Tool Settings 中配置 ModelSim，在 Simulation Sources 中确认顶层为 \code{tb_elinx}，执行 Run Behavioral Simulation。通过标志：末行含 \code{PASS checks=9238 errors=0}。

\textbf{纯 ModelSim 入口：}在 Transcript 中进入仓库根目录，执行以下命令。通过标志：7 行 \code{JUDGE CASE PASS} 和末行 \code{JUDGE ENTRY PASS cases=7}。汇总文件为 \code{out/modelsim/judge_*/result.txt}。
\begin{verbatim}
do sim/run_judge.do
\end{verbatim}

\textbf{完整 RTL 回归：}执行以下命令。通过标志：\RegCount{} 项均为 PASS。汇总写入 \code{out/regress/}，逐项输出写入 \code{out/sim/}。
\begin{small}
\begin{verbatim}
powershell -NoProfile -ExecutionPolicy Bypass `
  -File sim/regress.ps1
\end{verbatim}
\end{small}

\section{界面综合网表}
先在 eLinx 中执行 Run Synthesis 生成网表，再运行以下命令。第一条运行 NET-01；第二条先运行 RTL 五功能，再运行界面网表并比较 7 份端口记录（NET-02）。
\begin{small}
\begin{verbatim}
powershell -NoProfile -ExecutionPolicy Bypass `
  -File sim/run_gui_netlist.ps1

powershell -NoProfile -ExecutionPolicy Bypass `
  -File sim/run_gui_netlist.ps1 -Test tb_net_services
\end{verbatim}
\end{small}
通过标志：分别打印 \code{GUI NETLIST PASS tb_elinx} 和 \code{GUI NETLIST PASS tb_net_services}；自检 \ElinxChecks/0 和 \SvcChecks/0；五功能独立核对 65/0，7 份记录比较均一致。日志和汇总写入 \code{out/gui_netlist/}。

\section{命令行官方综合与网表验证}
执行官方综合引擎脚本，按工程配置直接生成 VQM 格式网表（耗时约 \CliSynthSec{}~s）。
\begin{small}
\begin{verbatim}
powershell -NoProfile -ExecutionPolicy Bypass `
  -File elinx/synth.ps1
\end{verbatim}
\end{small}

随后运行顶层五功能联合测试的网表仿真（NET-04），将生成的官方 VQM 网表输入给统一仿真入口：
\begin{small}
\begin{verbatim}
$net = (Get-Content out/latest_AFDX_top.json -Raw | ConvertFrom-Json).netlist
powershell -NoProfile -ExecutionPolicy Bypass `
  -File sim/run_gui_netlist.ps1 -Test tb_net_services -Netlist $net
\end{verbatim}
\end{small}

通过标志：自检打印 \code{[RESULT] tb_net_services PASS checks=1123 errors=0}，独立检查通过且 7 份端口记录与同激励 RTL 逐字节完全相同，Error、Fatal、Warning 均为 0。

\section{代码覆盖率复现}
在配置 ModelSim 的环境下，执行以下命令运行完整 38 项回归并开启 \code{sbceft} 代码覆盖率插桩统计：
\begin{small}
\begin{verbatim}
powershell -NoProfile -ExecutionPolicy Bypass `
  -File sim/regress.ps1 -Coverage
\end{verbatim}
\end{small}
覆盖率结果自动汇总于 \code{out/regress/coverage/}，基准存档于 \code{sim/reference/coverage/}。

\chapter{证据文件索引}\label{app:evidence}

\section{归档约定}
文件名除特别说明外均相对于 \code{docs/evidence/}。\code{manifest.json} 记录每个副本的名称（\code{file}）、原始位置（\code{source}）、字节数（\code{bytes}）和 SHA-256 摘要（\code{sha256}）。原始位置为生成时所在机器的绝对路径；复现输出以本次运行目录为准。回归、评委入口、性能、容量和异常注入的记录为 2026 年 10 月 14 日的结果，网表结果日期见第~\ref{ch:netlist}~章。开发期专项记录单独列出。

\section{标准测试与综合证据}
\begin{longtable}{@{}p{5.6cm}p{9.6cm}@{}}
\caption{标准测试及综合网表证据文件索引}\\
\longtablehead{2}{文件 & 对应测试与字段或日志定位}
\code{performance.json} & \code{tb_sys_perf}；\code{items/errors}，\code{reports.system.performance} 下各阶段、方向的吞吐、时延与抖动。 \\
\code{rtc.json} & \code{tb_sys_rtc}；\code{items/errors}、\code{idle_max_ns}、\code{loaded_max_ns}、\code{gbe_max_ns}、\code{gbe_mbps}。 \\
\code{capacity.log} & \code{tb_judge_capacity}；\code{[CAPACITY_SUMMARY]} 中发送数、接收配置数、抽测数与首末索引，\code{[RESULT]} 中检查项数和错误数。 \\
\code{snmp.log} & \code{tb_judge_snmp}；\code{[RESULT]} 的检查项数、错误数及 PASS 状态。 \\
\code{gui_behavior.log} & 界面行为仿真 \code{tb_elinx}；\code{[RESULT]} 和仿真时间，记录生成于 2026-10-12。 \\
\code{gui_elinx.log} & 界面网表 \code{tb_elinx}；\code{[RESULT]}。 \\
\code{gui_elinx_summary.json} & \code{results}、\code{diagnostics}、\code{passed}、\code{started/finished} 和 \code{jobs[].seconds}。 \\
\code{gui_services_summary.json} & 界面网表 \code{tb_net_services}；\code{results}、\code{check.items/errors}、\code{exact_external_records}、\code{adaptation}、\code{inputs}、\code{jobs}。 \\
\code{cli_services.log} & 命令行网表 \code{tb_net_services}；\code{[FUNCTION]}、\code{[RESULT]}。 \\
\code{cli_services_summary.json} & \code{checks/errors/result}、\code{check.items/errors}、\code{check.checks}、\code{check.download}、\code{inputs}、\code{jobs}。 \\
\code{cli_validation.json} & 命令行网表清单的 \code{net[]}：\code{test/tops/checks}、\code{independent_checks}、\code{same_external_records}；五功能见 \code{services_net}。 \\
\code{gui_resources.rpt} & 界面综合；\code{Total LUTs/FFs/CLBs/Memorys/Pins}。 \\
\code{gui_hierarchy.rpt} & 界面综合；\code{Resource Utilization by Entity} 表中的实体、逻辑单元、寄存器和存储位数。 \\
\code{device_capacity.rpt} & 器件资源容量；各资源行斜线后的总量用于计算占比，斜线前为该次布局结果。 \\
\code{gui_timing.json} & 界面综合时序只读查询；\code{rows[]} 中时钟 \code{clk} 的周期、setup、hold 与 Fmax；时间单位见 \code{time_unit}。 \\
\code{gui_synthesis_timing.pb} & 界面综合 SQLite 库；\code{elinx_intra_clock_summary} 及 \code{elinx_intra_clock_}\allowbreak\code{longest_sync_path_report}。 \\
\code{coverage/totals.json} & 代码覆盖率基准汇总（\code{sim/reference/coverage/}）；语句、分支、条件、状态机和翻转覆盖率及回归耗时。 \\
\code{coverage/modules.csv} & 21 个 RTL 模块的细分覆盖率明细表。 \\
\code{coverage/vcover_report.txt} & ModelSim \code{vcover} 导出的完整覆盖率文本报告。 \\
\code{out/latest_AFDX_top.json} & 命令行官方综合生成摘要；定位最新综合 VQM 网表、资源报告、时序数据库及 \code{tool_metrics.json}。 \\
\code{latest_tb_net_services_gui} & 命令行官方网表五功能仿真摘要；包含 \code{results}（1123/0）、\code{check}（65/0）、7 份一致记录及耗时 \NetCliTime{}~s。 \\
\end{longtable}

\section{异常注入证据}
\begin{longtable}{@{}p{5.6cm}p{9.6cm}@{}}
\caption{异常注入与计数读回证据文件索引}\\
\longtablehead{2}{文件 & 对应测试与字段或记录类型}
\code{fault_checks.json} & \code{tb_sys_fault}；\code{items/errors} 和 \code{reports} 中双端收发、GMII、系统阶段统计。 \\
\code{fault_events.txt} & 同测试的系统记录；\code{A/J/R/P/V} 行及 \code{G CLOCK_NS}，定位注入、阶段与配置口读回。 \\
\code{fault_rx_events.txt} & 同测试端 1 接收记录；\code{R/I} 行定位复位与就绪，\code{P/E} 行定位阶段与接收事件。 \\
\code{load_checks.json} & \code{tb_sys_load}；\code{items/errors}、\code{reports} 中各阶段的发送拒绝与交付核对。 \\
\code{load_events.txt} & 同测试系统记录；\code{V} 行为计数器读回，\code{P} 行为阶段。 \\
\code{load_tx_events.txt} & 同测试端 1 发送记录；\code{M/E} 行对应消息接纳和发送事件。 \\
\code{load_commands.txt} & 同测试输入向量；\code{M} 行给出提交消息，与发送事件和阶段对应。 \\
\code{overflow_checks.json} & \code{tb_sys_ovf}；\code{items/errors}、\code{reports} 中缓存溢出与恢复检查。 \\
\code{overflow_events.txt} & 同测试系统记录；\code{P/V} 行定位溢出阶段与计数读回。 \\
\code{overflow_tx_events.txt} & 同测试端 1 发送记录；\code{G/P/M/E} 行定位参数、阶段、消息与事件。 \\
\code{queue_overflow_checks.json} & \code{tb_rx_ovf_dq}；\code{items/errors}、\code{pulses}、\code{deliveries}、\code{stages.dq}。 \\
\code{queue_overflow_events.txt} & 同测试接收记录；\code{G RX_DQ_AW}、\code{P}、\code{E} 行定位队列参数、阶段和交付／溢出事件。 \\
\end{longtable}

\section{开发期专项记录}
\begin{longtable}{@{}p{5.6cm}p{9.6cm}@{}}
\caption{设计取舍与专项实验的归档索引}\\
\longtablehead{2}{文件 & 实验内容与字段}
\code{contention.json} & 同拍释放及服务争用；\code{runs[].measurements}，方向数据见其中的 \code{directions}。 \\
\code{route_default.json} & 默认布局布线实验；\code{started}、\code{settings}、\code{utilization}、\code{timing}。 \\
\code{route_timing.json} & 同实验的数据库查询；\code{table}、\code{time_unit}、\code{rows[]}。 \\
\code{vo_functional.json} & 整顶层 \code{.vo} 功能实验；\code{sdf}、\code{result_rows}、\code{first_failure}、\code{passed}。 \\
\code{vo_ram.json} & 不同位宽 RAM 对照；各项的 \code{width}、\code{rtl}、\code{vo_sim} 中检查数、错误数及失败记录。 \\
\code{vo_ram_bits.json} & RAM 位序分析；\code{width}、\code{input_bit_to_source_bit}、\code{output_source_index_identity}、\code{permuted}。 \\
\code{mac_results.json} & 自研 MAC 与官方配置对照；\code{baseline}、\code{idle_crc} 的 \code{latencies}，各工况的 \code{summary}；\code{estimates} 为估算字段。 \\
\code{mac_official_resources.json} & 官方配置资源；\code{utilization} 与 \code{packed_utilization} 分别保存综合和装箱报告。 \\
\code{mac_self_resources.json} & 自研 MAC 资源；同上两个字段，单口统计。 \\
\code{mac_official_stress.json} & 官方配置持续收发；\code{analysis.metrics/summary}、\code{jobs} 和 \code{tail}。 \\
\code{mac_self_rtl.json} & 自研 MAC 的 RTL 持续收发；\code{analysis.metrics/summary}、\code{jobs}。 \\
\code{mac_self_net.json} & 自研 MAC 的网表持续收发；\code{analysis.metrics/summary}、\code{jobs}。 \\
\code{mac_bad_crc.json} & 官方配置坏 FCS 实验；\code{tail} 中捕获记录及后续分析字段。 \\
\end{longtable}

\section{参考汇总与补充记录}
\code{sim/reference/} 下的两份汇总未列入 \code{manifest.json}。\code{regression.json}：\code{Started} 为开始时间，\code{Runs[]} 的 \code{Test/Param/Status/Checks/CheckItems} 对应表~\ref{tab:regress}，\code{Folder} 为各次运行输出目录。\code{judge_result.txt}：\code{JUDGE CASE} 为各用例的检查数、错误数、Warning 数和耗时，\code{JUDGE ENTRY} 为总结果，\code{JUDGE METRIC} 为性能、RTC 和突发指标。

\code{docs/evidence/version_match.json} 同样未列入 manifest，记录 \code{head}、\code{python_files}、\code{non_text_change} 和源码匹配统计，仅作版本核对的补充。

\section{数字来源汇总}
报告中的数字统一由 \code{docs/design/numbers.tex} 引入，来源见下表；测试配置与测量口径见第~\ref{ch:env}~章和第~\ref{ch:perf}~章。

\begin{longtable}{@{}p{5.6cm}p{9.6cm}@{}}
\caption{报告数字与原始来源对应}\\
\longtablehead{2}{数字或数字组 & 来源与取值位置}
\RegCount{} 项、\RegChecks{} 自检项、\RegIndep{} 独立项 & \code{sim/reference/regression.json} 的 \code{Runs[]} 长度、\code{Checks} 之和及非空 \code{CheckItems} 之和；逐项数值见完整回归表。 \\
评委入口 7 项、0 错误、0 Warning；各例耗时及总耗时 626.779 s & \code{sim/reference/judge_result.txt} 的 \code{JUDGE CASE}、\code{JUDGE ENTRY}，不是 README 中另一轮示例耗时。 \\
\ElinxChecks/0、\SvcChecks/0、65/0；\NetElinxTime{}、\NetGuiTime{}、\NetCliTime{} s & 界面摘要与最新命令行五功能摘要的 \code{results}（自检）和 \code{check.items/errors}（独立核对）；耗时见各摘要 \code{jobs} 中 \code{name=simulation} 的 \code{seconds}。 \\
界面 7 份、命令行 7 份一致记录；下载 41 项、16 字节 & 界面与命令行五功能摘要的 \code{exact_external_records}；两者的 \code{check.download.items} 及 \code{check.download.cases}。 \\
427486 个输入、3933 个进位参数 & \code{gui_services_summary.json} 的 \code{adaptation} 下 \code{added_input_defaults} 各计数之和，以及 \code{carry_use_parameters_added} 中的 \code{cin_used}。 \\
官方综合耗时 \CliSynthSec{} s & \code{out/synth/AFDX_top/.../tool_metrics.json} 的总运行耗时。 \\
网表日期与起止时间；CLB 4620 & 各摘要的 \code{started/finished}；整理后综合结果来自 \code{docs/FACTS.md} 第三节，当前 CLB 见 \code{gui_resources.rpt}。 \\
\CapChecks{}、\SnmpChecks{}；128 条发送、4096 条接收配置与抽测索引 & \code{capacity.log} 的 \code{[RESULT]}、\code{[CAPACITY_SUMMARY]} 及 \code{snmp.log} 的 \code{[RESULT]}。 \\
吞吐 \Thru{}、\ThruEightyOne{}、\ThruTen{} Mbps；时延 \Lat{}、\Ttx{}、\Trx{} µs；抖动 \Jit{} µs & \code{performance.json} 的 \code{reports.system.performance}，对应阶段、方向中的吞吐、\code{T_tx}、\code{T_rx}、\code{T_e2e_link_removed} 及抖动统计。 \\
突发清空时间 \BurstSevSeven{}、\BurstEightOne{} µs；RTC 误差 \RtcIdle{}、\RtcLoad{} ns & \code{judge_result.txt} 的 \code{JUDGE METRIC}；RTC 另见 \code{rtc.json} 的 \code{idle_max_ns}、\code{loaded_max_ns}、\code{gbe_max_ns}。 \\
资源、占比和综合时序数字 & \code{gui_resources.rpt}、\code{device_capacity.rpt}、\code{gui_timing.json} 与官方命令行综合查询结果。占比按综合用量除以器件总量计算。 \\
覆盖率 \CovStmt{}\% 语句、\CovBranch{}\% 分支、\CovCond{}\% 条件、\CovFsm{}\% 状态机、\CovToggle{}\% 翻转 & \code{sim/reference/coverage/totals.json} 的 \code{metrics} 及 \code{modules.csv} 各模块汇总。 \\
工具版本、依赖版本与五小时预算 & \code{README.md}“软件准备”、\code{FACTS.md} 第四节；VQM 适配依据 \code{tools/adapt_gui_vqm.py}。 \\
\end{longtable}
